\chapter{Metodología}
\label{cap:metodologia}

\section{Fuente de datos y cohorte}

La fuente primaria es la base abierta de COVID-19 de la Dirección General de
Epidemiología (DGE) de la Secretaría de Salud. Se emplea el corte anual
\texttt{\cohorteFuente}, que contiene \nBaseCompleta{} registros de notificación
del año 2021 con variables demográficas, entidad de residencia, comorbilidades y
desenlaces.

Como fuente de contexto se emplea el índice de marginación por entidad
federativa de CONAPO, edición 2020, construido sobre el Censo de Población y
Vivienda 2020 del INEGI~\parencite{conapo2020}. Se usa el archivo por entidad
(\texttt{IME\_2020.xls}), del que se toman el índice y el grado de marginación.

La cohorte de análisis se restringe a \textbf{casos COVID confirmados}
(\texttt{CLASIFICACION\_FINAL} $\in \{1,2,3\}$), lo que produce
$n=\cohorteN$ registros con una mortalidad observada de \cohorteTasaDef\,\%. La
restricción es deliberada: un modelo de mortalidad \emph{por COVID} debe
estimarse sobre enfermos de COVID. Incluir casos negativos y sospechosos mezcla
poblaciones con mecanismos de riesgo distintos y diluye la señal; la prevalencia
del desenlace cae entonces a \tasaDefBase\,\% y el modelo pasa a describir
mortalidad en población tamizada, que es otra pregunta.

\section{Definición de desenlaces y control de fuga de información}
\label{sec:fuga}

El desenlace modelado es la \textbf{defunción}, derivada del campo
\texttt{FECHA\_DEF}: la base codifica \texttt{9999-99-99} para paciente vivo.

\nota{La cohorte tiene una propiedad estructural que condiciona la lectura de
todo el capítulo de resultados y conviene declarar aquí. Entre los casos
confirmados, el número de defunciones registradas como ambulatorias es
\defAmbulatorias: fallecer implica estar registrado como hospitalizado. El
desenlace solo es posible, por tanto, en el \fracHosp\,\% de la cohorte
($n=\nHosp$), y el resto son casos que no pueden fallecer por construcción. Esto
no invalida la tarea --estimar riesgo al primer contacto sobre todo caso
confirmado es precisamente el triage que motiva la tesis-- pero significa que
parte de la discriminación consiste en separar ambulatorios de hospitalizados. La
sección~\ref{sec:sensibilidad} cuantifica cuánta.}

El control de fuga de información es la decisión metodológica más importante del
canal. El punto de predicción es el \emph{primer contacto}, de modo que ninguna variable
posterior a ese momento puede ser predictor. Se excluyen explícitamente:

\begin{itemize}[leftmargin=*]
  \item \texttt{UCI} e \texttt{INTUBADO}, que describen el curso hospitalario y
    solo se conocen después de la decisión que el modelo pretende apoyar.
  \item \texttt{FECHA\_DEF} y todo lo derivado de ella, salvo el desenlace mismo.
  \item \texttt{TIPO\_PACIENTE} como predictor de mortalidad, por ser posterior
    al triage.
\end{itemize}

Los predictores admitidos son la edad, el sexo y once comorbilidades o signos
registrados al primer contacto: diabetes, EPOC, asma, inmunosupresión,
hipertensión, otra comorbilidad, enfermedad cardiovascular, obesidad, enfermedad
renal crónica, tabaquismo y neumonía.

\nota{La inclusión de \texttt{NEUMONIA} merece justificación. En el formato de la
DGE se captura en el primer contacto y no durante la evolución hospitalaria, por
lo que es admisible en ese punto. Aun así resulta el predictor más informativo del
modelo (sección~\ref{sec:importancia}), de modo que la validez de esa captura es
un supuesto del que dependen los resultados. La
sección~\ref{sec:sensibilidad_res} reentrena el modelo sin ella para acotar
cuánto del desempeño y de la inequidad descansan en esa única variable.}

\section{Limpieza y recodificación}

La base codifica las categóricas con enteros y usa códigos especiales para
ausencia: \texttt{97} (no aplica), \texttt{98} (se ignora) y \texttt{99} (no
especificado). El tratamiento aplicado es:

\begin{itemize}[leftmargin=*]
  \item \textbf{Edad.} Valores fuera de $[0,120]$ se marcan como faltantes.
  \item \textbf{Sexo.} Se deriva un indicador binario \texttt{sexo\_mujer}
    ($1=$ mujer), dejando faltante el código \texttt{99}.
  \item \textbf{Entidad.} \texttt{ENTIDAD\_RES} se traduce a nombre mediante el
    catálogo de 32 entidades; los códigos 97--99 quedan sin etiqueta.
  \item \textbf{Comorbilidades.} Cada variable produce dos columnas: una binaria
    ($1=$ sí, $0=$ no, faltante en otro caso) y un \emph{indicador de faltante}.
    La codificación de ausencia se conserva como información en lugar de
    imputarse a ciegas, porque su frecuencia difiere entre entidades y es en sí
    un posible mecanismo de inequidad.
  \item \textbf{Grupo etario.} Cortes en 0--17, 18--39, 40--59 y 60 o más años,
    definidos una única vez para todo el canal.
\end{itemize}

\section{Partición de datos}
\label{sec:particion}

Se emplea una partición en tres bloques, estratificada por el desenlace con
semilla fija:

\begin{enumerate}[leftmargin=*]
  \item \textbf{Prueba} (25\,\%, $n=\testN$). Se separa primero y no interviene en
    ninguna decisión de ajuste.
  \item \textbf{Calibración} (20\,\% del resto, $n=\nCalib$). Se usa exclusivamente para
    ajustar la transformación de calibración y los umbrales por grupo de la
    mitigación post-procesamiento.
  \item \textbf{Entrenamiento} (el remanente, $n=\nTrain$). Ajuste de los modelos.
\end{enumerate}

El tercer bloque es necesario, y no un refinamiento cosmético. Tanto la
calibración como los umbrales por entidad son parámetros ajustados a datos; si se
estiman sobre el conjunto de prueba, las métricas resultantes están sesgadas
optimistamente. La sección~\ref{sec:mitigacion_res} cuantifica esa diferencia.

\section{Modelo de referencia}

Se entrenan dos modelos sobre los mismos predictores:

\begin{description}[leftmargin=*,style=nextline]
  \item[Regresión logística] Con imputación por mediana, estandarización y
    \texttt{class\_weight="balanced"}. Sirve como referencia interpretable.
  \item[Impulso del gradiente por histogramas] \texttt{HistGradientBoosting\-Classifier}
    con 300 iteraciones, tasa de aprendizaje 0.05, profundidad máxima 6 y
    \texttt{class\_weight="balanced"}. Es el modelo principal y el sujeto de la
    auditoría.
\end{description}

Ninguno de los dos incorpora la entidad federativa como predictor: la entidad se
reserva como atributo sensible para auditar, no como variable de entrada.

\subsection{Calibración}

El desbalance de clases del desenlace es de \desbalance{} a 1, y el parámetro
\texttt{class\_weight} en modo \texttt{balanced} lo compensa reponderando la
pérdida en proporción inversa a la frecuencia de cada clase. Como
consecuencia, el modelo sobreestima el riesgo absoluto por un factor cercano a
cuatro (sección~\ref{cap:resultados}). Se aplica por tanto una \textbf{regresión isotónica} ajustada sobre el
bloque de calibración. Al ser una transformación monótona, preserva el
ordenamiento --y por tanto el AUROC-- y corrige el nivel de la probabilidad.

Esta corrección es un requisito previo a la auditoría, no un adorno: sin ella, el
error de calibración medido por grupo queda dominado por el sesgo global de
entrenamiento y no por la inequidad entre grupos, que es el objeto de estudio.

\subsection{Punto de operación}

Todas las métricas dependientes de umbral se reportan en un punto de operación
único: el umbral $\tau$ que alcanza una \textbf{sensibilidad global de
\sensObjetivo\,\%}, propia de una herramienta de triage cuyo costo de omitir una
defunción excede con holgura el de una alarma innecesaria. Para el modelo
calibrado resulta $\tau=\umbralOperacion$. Fijar el mismo punto de operación en
todas las etapas es lo que hace comparables las tablas de desempeño, auditoría y
mitigación entre sí.

\section{Métricas}

\begin{table}[htbp]
\centering
\caption{Métricas empleadas por dimensión de evaluación.}
\label{tab:metricas}
\small
\begin{tabular}{p{3.4cm} p{10.4cm}}
\toprule
Dimensión & Métricas \\
\midrule
Discriminación & AUROC; AUPRC; sensibilidad y tasa de falsos positivos en el punto de operación \\
Calibración & Puntaje de Brier; curva de calibración por deciles; riesgo medio predicho frente a observado, global y por grupo \\
Equidad & Brecha máxima de TPR (\emph{equalized odds}); brecha de FPR; brecha de AUROC; brecha de tasa de selección (paridad demográfica); error máximo de calibración por grupo \\
Interpretabilidad & Importancia por permutación sobre el AUROC \\
Utilidad agregada & Exactitud balanceada en el punto de operación \\
\bottomrule
\end{tabular}
\end{table}

Las métricas de equidad se implementan explícitamente en lugar de delegarse a una
biblioteca, para que el cálculo sea auditable línea por línea; son equivalentes a
las definiciones de Fairlearn~\parencite{fairlearn2020}.

\subsection{Criterio de soporte por grupo}

Un grupo se considera evaluable si tiene al menos 500 observaciones y 20 casos
positivos en el bloque correspondiente. Por debajo de ese soporte, la estimación
de sensibilidad tiene una varianza que domina cualquier señal de inequidad. Los
grupos descartados se reportan explícitamente en lugar de omitirse en silencio;
con la cohorte empleada, las \nEntidades{} entidades superan el criterio.

En el análisis interseccional el mínimo se relaja a \interMinN{} observaciones.
Exigir \numprint{500} dentro de un subgrupo etario dejaría fuera precisamente a
las entidades pequeñas, que son las que interesa examinar. La
sección~\ref{sec:interseccional} reporta el resultado bajo ambos criterios para
que la conclusión no dependa de esa elección.

\section{Estrategias de mitigación}

Se evalúan las tres familias bajo el mismo punto de operación:

\begin{description}[leftmargin=*,style=nextline]
  \item[Pre-procesamiento] Reponderación de muestras inversamente proporcional al
    número de registros de cada entidad, de modo que los estados con menos datos
    pesen más durante el entrenamiento. Se reentrena el modelo con los
    hiperparámetros idénticos a los de la línea de base, para que la única
    diferencia sean los pesos.
  \item[In-procesamiento] Reducción por gradiente exponenciado con restricción de
    \emph{equalized odds}~\parencite{fairlearn2020}, sobre una submuestra de
    \numprint{120000} registros del bloque de entrenamiento y con un clasificador
    lineal como aprendiz base, dado que el método reentrena el modelo en cada
    iteración. Esta limitación la hace no comparable uno a uno con las demás
    intervenciones, y así se reporta.
  \item[Post-procesamiento] Un umbral por entidad, elegido para alcanzar la
    sensibilidad objetivo \emph{dentro} de cada estado. Los umbrales se estiman en
    el bloque de calibración y se evalúan en el de prueba.
\end{description}

\subsection{Frontera de Pareto}

La frontera se traza interpolando linealmente entre el umbral global y el umbral
por entidad,
\begin{equation}
  \tau_g(\alpha) = (1-\alpha)\,\tau + \alpha\,\tau_g,
  \qquad \alpha \in [0,1],
\end{equation}
donde $\tau_g$ es el umbral estimado para la entidad $g$. Cada valor de $\alpha$
define una política de despliegue: $\alpha=0$ es el modelo sin mitigar y
$\alpha=1$ la corrección completa por entidad. El conjunto de pares
(brecha de sensibilidad, desempeño global) constituye el instrumento de decisión
entregable, porque traslada al responsable de la política una elección que es
normativa y no técnica.

\section{Análisis de sensibilidad: subconjunto hospitalizado}
\label{sec:sensibilidad}

Se reentrena el modelo, con los mismos hiperparámetros y el mismo esquema de
partición, restringido a los pacientes hospitalizados. Es la población en la que
el desenlace es posible, la que corresponde a los modelos clínicos de mortalidad
hospitalaria con los que suele compararse un trabajo de este tipo, y la que
tendría una eventual validación externa sobre bases de egresos.

Se reentrena en lugar de restringir a posteriori el conjunto de prueba del modelo
global, porque ese modelo habría aprendido de ambulatorios y la comparación
mediría otra cosa. El AUPRC se reporta junto con su razón respecto de la
prevalencia, que es su línea base: sin normalizar, dos AUPRC de cohortes con
prevalencias de 5.6\,\% y 47.5\,\% no son comparables.

\section{Cuantificación de la incertidumbre}
\label{sec:incertidumbre}

Las brechas se calculan como diferencia entre el máximo y el mínimo de 32 grupos
de tamaños muy desiguales. Ese estadístico está sesgado al alza: aun sin ningún
efecto de entidad, el máximo y el mínimo de 32 estimaciones ruidosas se separan
por azar. Reportarlo como estimación puntual, sin referencia, invita a
sobreinterpretarlo. Se emplean por ello dos procedimientos complementarios.

El primero es un \textbf{remuestreo estratificado por entidad}: \boB{} réplicas
en las que cada entidad se remuestrea con reemplazo conservando su tamaño, de las
que se obtienen intervalos de confianza al 95\,\% por percentiles para el AUROC y
la sensibilidad de cada estado. Se estratifica dentro de la entidad, y no sobre
todo el conjunto, para que el intervalo refleje la precisión con que se estima
\emph{ese} estado; las entidades pequeñas deben salir con intervalos anchos.

El segundo es una \textbf{prueba de permutación}: \boP{} reasignaciones al azar
de la etiqueta de entidad, conservando los tamaños de grupo, que construyen la
distribución de la brecha bajo la hipótesis nula de que la entidad no influye en
el desempeño. El valor $p$ se calcula con la corrección de continuidad habitual,
$(\#\{\text{nula} \ge \text{observada}\}+1)/(P+1)$, de modo que nunca resulta
exactamente cero. Es este contraste, y no el intervalo de confianza, el que
permite afirmar si las disparidades son \emph{estadísticamente significativas},
que es la formulación de la hipótesis~\ref{h1}.

\section{Transportabilidad geográfica interna}
\label{sec:loeo}

La partición aleatoria reparte cada entidad entre entrenamiento y prueba, de modo
que el modelo siempre ha visto casos del estado en el que se le evalúa. Eso mide
interpolación dentro de una población conocida, no generalización a un contexto
nuevo, que es lo que importa para desplegar en una entidad cuyo perfil no estaba
en los datos de ajuste.

Se entrenan por ello \loeoNent{} modelos: cada uno excluye por completo una
entidad y se evalúa sobre la totalidad de sus registros. La calibración se ajusta
sobre un bloque tallado de las entidades restantes, y el umbral se fija para
alcanzar la sensibilidad objetivo \emph{en esas entidades} --nunca en la
excluida--, porque en despliegue no se conoce el desenlace del estado nuevo. La
pérdida por transportabilidad es la diferencia entre el AUROC de una entidad
cuando participó en el entrenamiento y cuando no.

\nota{Este diseño \emph{no} pone a prueba la hipótesis~\ref{h4}, que habla de una
fuente independiente: aquí la fuente, el año y el proceso de captura son los
mismos. Lo que somete a prueba es la generalización geográfica, que es el eje del
trabajo. \ref{h4} sigue pendiente (sección~\ref{sec:trabajo_futuro}).}

\section{Reproducibilidad}

El canal se implementa en cuatro cuadernos ejecutados en orden, con semillas
fijas y un contrato de datos explícito entre etapas que verifica procedencia,
esquema y cohorte antes de consumir cualquier artefacto. Las cifras del
capítulo~\ref{cap:resultados} no se transcriben a mano: se generan a partir de los
archivos de resultados que exportan los cuadernos. El apéndice~\ref{ap:reproducibilidad}
documenta el mecanismo.
